% Preámbulo compartido — Informe Técnico Scapegoat Tree
\usepackage[utf8]{inputenc}
\usepackage[T1]{fontenc}
\usepackage[spanish,es-tabla]{babel}
\usepackage{mathptmx}
\usepackage{textcomp}
\usepackage{microtype}
\usepackage{geometry}
\usepackage[dvipsnames]{xcolor}
\usepackage{graphicx}
\usepackage{booktabs}
\usepackage{longtable}
\usepackage{array}
\usepackage{multirow}
\usepackage{caption}
\usepackage{subcaption}
\usepackage{float}
\usepackage{amsmath,amssymb,amsthm}
\usepackage{mathtools}
\usepackage{listings}
\usepackage{fancyhdr}
\usepackage{titlesec}
\usepackage{enumitem}
\usepackage{tikz}
\usepackage{pgfplots}
\usetikzlibrary{arrows.meta,positioning,shapes.geometric,fit,backgrounds,calc}

\pgfplotsset{compat=1.18}

\geometry{left=3cm, right=2.5cm, top=2.8cm, bottom=2.8cm}
\setlength{\parskip}{0.45em}
\setlength{\parindent}{1.2em}
\linespread{1.15}

\definecolor{codebg}{RGB}{244,246,248}
\definecolor{codeframe}{RGB}{49,87,213}
\definecolor{esanblue}{RGB}{26,26,46}

% Lenguaje Go para listings (no viene por defecto)
\lstdefinelanguage{Go}{
  morekeywords={break,case,chan,const,continue,default,defer,else,fallthrough,for,func,go,goto,if,import,interface,map,package,range,return,select,struct,switch,type,var,bool,byte,complex64,complex128,error,float32,float64,int,int8,int16,int32,int64,rune,string,uint,uint8,uint16,uint32,uint64,uintptr,true,false,nil},
  sensitive=true,
  morecomment=[l]{//},
  morecomment=[s]{/*}{*/},
  morestring=[b]",
  morestring=[b]',
}

\lstdefinestyle{goi}{
  basicstyle=\ttfamily\footnotesize,
  breaklines=false,
  columns=fullflexible,
  keepspaces=true,
  showstringspaces=false,
}

% Código inline Go con caracteres especiales (genéricos, llaves, etc.)
\newcommand{\go}[1]{\lstinline[style=goi]|#1|}

\lstdefinestyle{gocode}{
  language=Go,
  basicstyle=\ttfamily\footnotesize,
  keywordstyle=\color{blue}\bfseries,
  commentstyle=\color{gray}\itshape,
  stringstyle=\color{teal},
  numbers=left,
  numberstyle=\tiny\color{gray},
  stepnumber=1,
  numbersep=8pt,
  backgroundcolor=\color{codebg},
  frame=single,
  rulecolor=\color{codeframe},
  breaklines=true,
  tabsize=4,
  showstringspaces=false,
  literate={á}{{\'a}}1 {é}{{\'e}}1 {í}{{\'i}}1 {ó}{{\'o}}1 {ú}{{\'u}}1
           {ñ}{{\~n}}1 {Á}{{\'A}}1 {É}{{\'E}}1 {Í}{{\'I}}1 {Ó}{{\'O}}1 {Ú}{{\'U}}1 {Ñ}{{\~N}}1
}

\lstdefinestyle{sqlcode}{
  language=SQL,
  basicstyle=\ttfamily\footnotesize,
  backgroundcolor=\color{codebg},
  frame=single,
  rulecolor=\color{codeframe},
  breaklines=true
}

\lstdefinestyle{jsoncode}{
  basicstyle=\ttfamily\footnotesize,
  backgroundcolor=\color{codebg},
  frame=single,
  rulecolor=\color{codeframe},
  breaklines=true
}

\lstdefinestyle{pscode}{
  basicstyle=\ttfamily\footnotesize,
  backgroundcolor=\color{codebg},
  frame=single,
  rulecolor=\color{codeframe},
  breaklines=true,
  literate={\$}{{\$}}1
}

\captionsetup{
  font=small,
  labelfont=bf,
  justification=centering,
  skip=8pt
}

\captionsetup[lstlisting]{singlelinecheck=false}

\setlength{\headheight}{15pt}
\pagestyle{fancy}
\fancyhf{}
\fancyhead[L]{\small\textit{Algoritmos y Estructura de Datos --- ESAN 2026-1}}
\fancyhead[R]{\small\textit{Scapegoat Tree}}
\fancyfoot[C]{\thepage}
\renewcommand{\headrulewidth}{0.4pt}

\titleformat{\chapter}[display]
  {\normalfont\bfseries\color{esanblue}}
  {\chaptertitlename\ \thechapter}{12pt}{\Large}
\titlespacing*{\chapter}{0pt}{-10pt}{20pt}

\titleformat{\section}{\normalfont\bfseries\color{esanblue}}{\thesection}{0.8em}{}
\titleformat{\subsection}{\normalfont\bfseries}{\thesubsection}{0.8em}{}

\newtheorem{teorema}{Teorema}[chapter]
\newtheorem{invariante}[teorema]{Invariante}

\newcommand{\bigO}{\mathcal{O}}
\newcommand{\Scapegoat}{\textit{Scapegoat Tree}}

% hyperref al final del preámbulo
\usepackage{hyperref}
\usepackage{bookmark}

\hypersetup{
  colorlinks=true,
  linkcolor=NavyBlue,
  citecolor=NavyBlue,
  urlcolor=NavyBlue,
  pdftitle={Informe Tecnico --- Scapegoat Tree},
  pdfauthor={Universidad ESAN}
}
